% Please do not change %
\documentclass[11pt]{article} % font size
\usepackage[margin=1in]{geometry} % margins
\usepackage{amsmath,amsthm,amssymb} % for the  common "math symbols"
\usepackage{algorithm}
\usepackage{algpseudocode}
\usepackage[shortlabels]{enumitem} % for enumeration
\usepackage{booktabs} % if you need tables
\usepackage{listings}
\usepackage{color}
\DeclareMathOperator*{\argmax}{argmax}
\DeclareMathOperator*{\argmin}{argmin}

% \theoremstyle{definition}
\newtheorem{theorem}{Theorem}[section]
\newtheorem{corollary}[theorem]{Corollary}
\newtheorem{lemma}[theorem]{Lemma}
\newtheorem{proposition}[theorem]{Proposition}
\newtheorem{conjecture}[theorem]{Conjecture}
\newtheorem{definition}[theorem]{Definition}
\newtheorem{fact}[theorem]{Fact}
\theoremstyle{definition}
\newtheorem{example}[theorem]{Example}

\lstset{
  frame=none,
  xleftmargin=2pt,
  stepnumber=1,
  numbers=left,
  numbersep=5pt,
  numberstyle=\ttfamily\tiny\color[gray]{0.3},
  belowcaptionskip=\bigskipamount,
  captionpos=b,
  escapeinside={*'}{'*},
  language=haskell,
  tabsize=2,
  emphstyle={\bf},
  commentstyle=\it,
  stringstyle=\mdseries\rmfamily,
  showspaces=false,
  keywordstyle=\bfseries\rmfamily,
  columns=flexible,
  basicstyle=\small\sffamily,
  showstringspaces=false,
  morecomment=[l]\%,
}
% Feel free to include additional packages (if needed), but make sure it does not override the margin/font requirements.
%
\setlength{\parindent}{0pt} % no indent for new paragraphs.
\setlength{\parskip}{5pt} % distance between paragraphs, which will be used when you have a blank line in your LaTeX code.


%%%%%%%%%%%%%%%%%%%%%%%%%%%%%%%%%%%%%%%%%%%%%%%%%%%%%%%%%%%%%%%%%%%%%%%%%%%%%%
\title{Categories, Proofs, and Processes: \\Assignment 2}
%
\author{Wira Azmoon Ahmad}
\date{4 Feb 2024}
%
\begin{document}
%
\maketitle

%You may use the usual \LaTeX\ environments to structure your answers:
%\begin{align}
%        \label{basegnn}
%        \mathbf{h}_u^{(t+1)} = \sigma \bigg( \mathbf{W}_\text{self}^{(t)} \mathbf{h}_u^{(t)}  + \mathbf{W}_\text{neigh}^{(t)} \sum_{v \in N(u)} \mathbf{h}_v^{(t)} + \mathbf{b}^{(t)} \bigg): 
%\end{align}


%\begin{table}[h]
%\caption{A simple table.}
%\centering
%\begin{tabular}[t]{lrrrr}
%\toprule
%Graphs & Can & Be  & Fun & Too \\ 
%\midrule 
%$G_1$ & $1$  & $2$ & $3$ & $4$ \\ 
%$G_2$ & $-1$  & $-2$ & $-3$ & $-4$ \\ 
%$G_3$ & $0.1$  & $0.2$ & $0.3$ & $0.4$ \\ 
%\bottomrule
%\end{tabular}
%\label{tab}
%\end{table}

\section*{Question 1}
\addtocounter{section}{1}

\begin{proof}
\begin{equation}
\begin{aligned}
    (g_1 \times g_2)\circ (f_1 \times f_2) 
    &= (g_1 \times g_2)\circ \langle f_1 \circ \pi_1, f_2 \circ \pi_2 \rangle\\
    &= \langle g_1 \circ f_1 \circ \pi_1, g_2 \circ f_2 \circ \pi_2 \rangle\\
    &= (g_1 \circ f_1) \times (g_2 \circ f_2)\\
\end{aligned}
\end{equation}
\begin{equation}
\begin{aligned}
    \text{id}_{A_1}\times \text{id}_{A_2} 
    &= \langle \text{id}_{A_1} \circ \pi_1, \text{id}_{A_2} \circ \pi_2 \rangle\\
    &= \langle \pi_1 \circ \text{id}_{A_1 \times A_2}, \pi_2 \circ \text{id}_{A_1 \times A_2} \rangle\\
    &= \text{id}_{A_1 \times A_2}
\end{aligned}
\end{equation}
\end{proof}

\section*{Question 2}
\addtocounter{section}{1}



\end{document}
